Spiral galaxies are disk-like rotating structures, whose dynamical state is
fairly grasped by the so-called rotation curves, quantifying the mean
rotational velocity of the disk at different distances from the center (see
Figure 1, curve B). An interesting feature is the flat region of the curve,
which is attributed to the presence of dark matter. Without it, rotation
velocities would drop steadily at large radii from the center, as depicted in
curve A.

\ioaafig{2021_DA_DA1-f5.pdf}
% vs radius: curve A falling at large radii, curve B flattening (source paper
% Q1-1, "Figure 1: Rotation curves")

In disk galaxies a strong correlation has been observed between the intrinsic
luminosity of the whole galaxy and the asymptotic rotational velocity (as
given by the rotation curve for the outer edge of the galaxy, i.e.\
$R_{\max}$), a result that is known as the Tully--Fisher relation. This
relation also holds if you use the luminosity in a specific band. This is
shown on Figure 2 for a number of galaxies in a galaxy cluster. Every dot is a
galaxy, and the solid line is the best-fit linear relation between absolute
magnitude in $K$ band and $\log_{10}(V_{\max})$ for the whole sample.

\ioaafig{2021_DA_DA1-f1.pdf}

{\centering\small Figure 2: Absolute magnitude in $K$ band vs
$\log_{10}(V_{\max}\,[\mathrm{km\,s^{-1}}])$. Tully--Fisher relation for
several galaxies. Every dot represents a galaxy. The dark points are five
selected galaxies, for which we will provide some numbers in part 1.2.\par}

Another interesting trend is shown in Figure 3: disks with larger stellar
masses ($M_*$) tend to have smaller gas fractions ($M_{gas}/M_*$).

\ioaafig{2021_DA_DA1-f2.pdf}

In the following questions you will be asked to extract physical information
about the galaxies using the scaling relations just introduced. Consider the
following guidelines:
\begin{itemize}
\item Assume that $V_{\max}$ was measured at the same radius for all galaxies
  ($R_{\max}$), in the flat part of the rotation curves and well beyond the
  end of the stellar disk.
\item Use $M_{dm}$ for the dark matter mass up to $R_{\max}$ and $M_{tot}$
  for the sum of all components (gas, stars and dark matter).
\item Assume that all galaxies have identical stellar populations (the term
  stellar population refers to the type of stars that are present in a
  galaxy, and the relative amount of each type with respect to the total
  number of stars), and assume that the gaseous component does not interact
  with the stellar light.
\item The galaxy cluster is far away. Its distance is much larger than the
  cluster size.
\item In spherically-symmetric mass distributions, to infer the gravitational
  effect on a particle at distance $r$ from the center, it suffices to
  consider the total mass enclosed up to that radius $M(\leq r)$ as if it
  were placed at the very center of the distribution.
\end{itemize}

\medskip
\textbf{Part 1 (20 points).}

\begin{description}
\item[1.1] From an analysis of Figure 3, find the appropriate constants in the
  following relation: $M_{gas} = a \times M_*^b$
  \[ a = ?\qquad b = ? \]
  \hfill[5.0\,pt]
\item[1.2] In the plot of the Tully--Fisher relation there are 5 highlighted
  points. Data for these 5 galaxies is given in the following table. Use this
  dataset to find the appropriate constants for the TF relation presented
  below the table, by means of a linear fit using the method of least
  squares.

  \emph{Note:} Treat $\log_{10}(V_{\max})$ as the $x$ variable and $K$ as the
  $y$ variable in the linear fit.

  \begin{center}
  \begin{tabular}{cc}
  $V_{\max}$ [km/s] & $K$ [mag] \\
  \hline
   79.4 & $-16.8$ \\
  100.1 & $-19.2$ \\
  158.5 & $-21.3$ \\
  251.2 & $-21.4$ \\
  316.2 & $-24.0$ \\
  \end{tabular}
  \end{center}

  \[ K = c \times \log_{10}(V_{\max}) + d \]
  \[ c = ?\qquad d = ? \]
  \hfill[15.0\,pt]
\end{description}

\medskip
\textbf{Part 2 (16 points).}

For two galaxies, G1 and G2, in the cluster, the recorded \emph{apparent}
magnitudes are:
\[ k_1 = 19.2\,;\qquad k_2 = 25.2 \]
Using this information and the relations calibrated in Part 1, find the
correct exponents in the following equations:

\begin{description}
\item[2.1] $\dfrac{M_{*1}}{M_{*2}} = 10^e$;\quad $e = ?$ \hfill[6.0\,pt]
\item[2.2] $\dfrac{M_{gas1}}{M_{gas2}} = 10^f$;\quad $f = ?$ \hfill[4.0\,pt]
\item[2.3] $\dfrac{M_{tot1}}{M_{tot2}} = 10^g$;\quad $g = ?$ \hfill[6.0\,pt]
\end{description}

\medskip
\textbf{Part 3 (15 points).}

\begin{description}
\item[3.1] Fill in the missing values in the table using the fact that for
  galaxy $G_1$, the dark-to-baryonic mass ratio up to $R_{\max}$ is 6.82.
  \hfill[15.0\,pt]

  \begin{center}
  \begin{tabular}{cccccc}
  Galaxy & Apparent magnitude $k$ & $M_{gas}\,[M_\odot]$ & $M_*\,[M_\odot]$ &
  $M_{dm}\,[M_\odot]$ & $M_{tot}\,[M_\odot]$ \\
  \hline
  $G_1$ & 19.2 & & & & $4.39 \times 10^{11}$ \\
  \end{tabular}
  \end{center}
\end{description}

\medskip
\textbf{Part 4 (24 points).}

\begin{description}
\item[4.1] Consider a systematic uncertainty of $\sigma_{sys} = \pm 0.2$ in
  each apparent magnitude due to CCD calibration errors. Then $k_1$ must be
  read as $k_1 = 19.2 \pm 0.2$, i.e., the only thing we know is that $k_1$
  most likely lies in the interval $[19.0, 19.4]$. The same goes for $k_2$.

  Recalculate the exponent in the scaling relation
  $\frac{M_{*1}}{M_{*2}} = 10^e$ (found in 2.1), expressing $e$ as an
  interval estimated by considering the extreme possible variations in $k_1$
  and $k_2$.
  \[ e \in [?, ?] \]
  \hfill[4.0\,pt]
\item[4.2] Now we consider that there is always a natural spread of the data
  around any relation. For instance, for a given value of the $K$ magnitude
  the TF relation gives a single value of $\log_{10}(V_{\max})$, but it would
  be more realistic to report an interval of plausible values, derived from
  the natural spread of the data around the mean TF relation. We call this
  the statistical uncertainty, $\sigma_{stat}$.

  Estimate the statistical uncertainty if $\log_{10}(V_{\max})$ is inferred
  from $K$ using the TF relation from question 1.2. For this, consider for
  each point the difference between the value of $\log_{10}(V_{\max})$
  estimated from $K$ using your linear fit and the actual measurement of
  $\log_{10}(V_{\max})$, and take $\sigma_{stat}$ as two times the root mean
  square (RMS) of these differences (the RMS of a set of values is the square
  root of the arithmetic mean of the squares of those values).
  \[ \sigma_{stat} = ? \]
  \hfill[10.0\,pt]
\item[4.3] Recalculate the exponent in the scaling relation
  $\frac{M_{tot1}}{M_{tot2}} = 10^g$, expressing $g$ as an interval estimated
  by considering the extreme possible variations arising from both the
  systematic and statistical uncertainties:
  \[ g \in [?, ?] \]
  \hfill[10.0\,pt]
\end{description}
